% Created 2026-08-09 Sun 15:57
% Intended LaTeX compiler: pdflatex
% SUPERSEDED — see docs/006-deep-research-spec.md (authoritative research feature doc).
\documentclass[11pt]{article}

\usepackage{inputenc}
\usepackage{fontenc}
\usepackage{graphicx}
\usepackage{longtable}
\usepackage{url}
\usepackage{amsmath}
\usepackage{hyperref}
\date{\today}
\title{Early Research for a High-Availability Web Search Gateway\\\medskip
\large Notes on free capacity, provider limits, SearXNG, policy constraints, and a proposed Pi-controlled architecture}
\hypersetup{
 pdfauthor={},
 pdftitle={Early Research for a High-Availability Web Search Gateway},
 pdfkeywords={},
 pdfsubject={},
 pdfcreator={},
 pdflang={English}}

% minimal-texlive fix: itemize bullets + section sign need textcomp (tcrm), not installed
\renewcommand{\labelitemi}{$\bullet$}
\renewcommand{\labelitemii}{$\circ$}
\renewcommand{\labelitemiii}{$\cdot$}
\renewcommand{\labelitemiv}{$\bullet$}
\pdfstringdefDisableCommands{\def\textsubscript#1{#1}\def\textsuperscript#1{#1}}

\begin{document}

\maketitle
\setcounter{tocdepth}{3}
\tableofcontents

\section*{Status}
\label{sec:org4750991}

These are working notes from the early research and brainstorming phase.  They
record the evidence, calculations, rejected ideas, and current design direction;
they are not yet an approved implementation specification.

No code was changed during this research phase.

Primary local inputs:

\begin{itemize}
\item Existing extension: \texttt{extensions/web-search/}.
\item Existing skill: \texttt{skills/web-search/SKILL.md}.
\item User-provided report: \texttt{/home/pirackr/Working/generated-reports/web-search.pdf}.
\item Extracted report text used during review: \texttt{/tmp/web-search-report.txt}.
\item Exa terms text captured during review: \texttt{/tmp/exa-terms.txt}.
\end{itemize}
\section*{Executive conclusions}
\label{sec:orgfb06cb1}

\begin{itemize}
\item The desired scale is approximately \textbf{5,000 programmatic searches/day}, or
approximately \textbf{150,000/month}.
\item The average request rate is small: \texttt{5,000 / 86,400 = 0.0579 requests/second}.
Per-minute and per-second limits matter for report bursts, but \textbf{monthly free
credits are the binding constraint}.
\item Approximately \textbf{114 fresh searches/day} are available from the verified
recurring direct/AI-search free tiers (Exa, Brave, and Tavily basic).
\item Including fixed-quota third-party SERP intermediaries, especially Bright
Data, raises the documented recurring total to approximately \textbf{286/day}.
\item A broad optimistic ceiling, including variable generic scraping credits, is
approximately \textbf{330/day}, but it is not reliable planning capacity.
\item One-time signup and trial grants provide approximately \textbf{25,000--31,000
searches}, enough for only approximately \textbf{5--6 days} at the target rate.
\item Self-hosted SearXNG can accept 5,000 local calls/day, but it owns no search
index.  It fans queries out to upstream engines that can return CAPTCHAs,
throttle, block the host IP, or change markup.  It therefore cannot guarantee
fresh results or a 99\% fresh-search SLA.
\item A hard-\char36{}0 system cannot guarantee 99\% fresh, unrestricted, general-web search
at 5,000/day.  It can target \textbf{99\% non-empty tool responses} only because the
accepted success definition includes fresh cache hits, stale cache hits, and
explicit partial results.
\item The recommended baseline is a local resilient search gateway:
\texttt{fresh cache -> healthy authorized/live engines -> stale cache -> explicit
  partial response}.
\item Pi must continue to orchestrate every query and receive every result.  Hosted
consumer research agents are not substitutes for a raw search backend.
\item Quota routing must use one legitimate account per provider and stay within
every provider's authorization.  Multi-account quota multiplication, stealth
automation, proxy rotation, CAPTCHA solving, and similar circumvention are
excluded.
\end{itemize}
\section*{Mission and constraints}
\label{sec:orgc0adea0}

\subsection*{Workload}
\label{sec:org60c92c4}

\begin{itemize}
\item Approximately 1--10 reports/day.
\item A report may fan out to approximately 1,000 searches.
\item Peak daily workload may reach approximately 5,000 searches.
\item Preserve high fan-out; do not redefine the problem by forcing reports to use
only a small number of searches.
\item Expected cache reuse is only approximately 10--40\%.
\end{itemize}
\subsection*{Hard requirements}
\label{sec:orgb3a22c0}

\begin{itemize}
\item Maximum monthly operating cost: \textbf{\char36{}0}, including APIs, proxies, and servers.
\item Self-hosting is allowed when it uses already-available hardware/connectivity.
\item Pi controls query generation, dispatch, result receipt, synthesis, and report
production.
\item General-web queries must not be restricted to a small curated corpus.
\item Cached, stale, or partial non-empty results count toward the availability
objective, but their status must be explicit.
\item Fresh-only 99\% coverage is not required after clarification.
\item Respect provider terms, hard quotas, burst limits, \texttt{Retry-After}, and blocks.
\end{itemize}
\subsection*{Explicit non-goals}
\label{sec:org34bbcd6}

\begin{itemize}
\item No multiple accounts at one provider to multiply a free allowance.
\item No rotating identities, IPs, or proxies to evade a rate limit or block.
\item No browser fingerprint mimicry, stealth scraping, or CAPTCHA solving.
\item No automation of consumer chat subscriptions as unofficial search APIs.
\item No hidden conversion of a stale or partial result into a result labeled
\texttt{fresh}.
\item No claim that best-effort upstream scraping provides an SLA.
\end{itemize}
\section*{Existing extension architecture}
\label{sec:org2b4b739}

The current extension is intentionally small and provider-oriented:

\begin{itemize}
\item \texttt{extensions/web-search/index.ts} registers \texttt{web\_lookup} and \texttt{fetch\_web}.
\item \texttt{extensions/web-search/search.ts} implements \texttt{webLookup()} and walks an
ordered provider chain.  The first available provider returning results wins.
\item \texttt{extensions/web-search/types.ts} contains shared search/fetch result types.
\item \texttt{extensions/web-search/engines/exa.ts} calls Exa directly.
\item \texttt{extensions/web-search/engines/duckduckgo.ts} scrapes DuckDuckGo's HTML
endpoint and decodes redirect URLs.
\item \texttt{extensions/web-search/engines/tavily.ts} calls Tavily directly.
\item \texttt{extensions/web-search/strategies/readability.ts} fetches pages and extracts
readable content with Mozilla Readability and linkedom.
\item \texttt{tests/web-search.test.ts} contains the web-search extension tests.
\item \texttt{skills/web-search/SKILL.md} provides engine-selection judgment,
search-vs-fetch guidance, prompt-injection safety, and recovery guidance.
\end{itemize}

Current automatic chain:

\begin{enumerate}
\item Exa.
\item DuckDuckGo fallback.
\end{enumerate}

Tavily is registered but opt-in rather than part of \texttt{engine: "auto"}.

The extension already has process-level usage caps:

\begin{itemize}
\item \texttt{-{}-{}web-search-max-lookups}
\item \texttt{-{}-{}web-search-max-fetches}
\end{itemize}

Those caps prevent accidental use inside one process; they are not persistent
cross-session provider quota accounting and do not provide caching, health
tracking, or circuit breaking.

Historical compatibility note: the tool is named \texttt{web\_lookup} rather than
\texttt{web\_search}.  \texttt{web\_search} collides with Anthropic's reserved server-side web
search tool name and previously broke requests through an Anthropic-compatible
non-Anthropic gateway.
\section*{Capacity arithmetic}
\label{sec:orgd420882}

\subsection*{Baseline volume}
\label{sec:orgfb27f55}

\begin{verbatim}
5,000 searches/day * 30 days = 150,000 searches/month
5,000 searches/day / 86,400 seconds = 0.0579 searches/second average
1,000-search report completed in one hour = 16.7 searches/minute
\end{verbatim}

Consequently, most published burst limits are generous enough if requests are
queued.  The scarce resource is the monthly allowance.
\subsection*{Verified fixed recurring free capacity}
\label{sec:org41f4cf1}

The following calculation assumes one legitimate account per provider and
Tavily \texttt{basic} searches.

\begin{center}
\begin{tabular}{lllrll}
Provider & Recurring allowance & Approx. searches/month & 30-day average/day & Published burst/operational limit & Important caveat\\
\hline
Bright Data SERP API & 5,000 credits/month & 5,000 & 166.7 & No clear binding limit found & Third-party SERP scraping/unlocking, not an independent index\\
Exa & \char36{}10 recurring credit/month; basic search \char36{}7/1,000 & 1,428.6 & 47.6 & 10 QPS found in docs & Credit-denominated; \char36{}20 signup credit is separate\\
Brave Search API & \char36{}5 recurring credit/month; \char36{}5/1,000 & 1,000 & 33.3 & 50 RPS & Credit card required for free plan; multiple-account bypass prohibited\\
Tavily basic & 1,000 credits/month; 1 credit/basic query & 1,000 & 33.3 & 100 RPM development limit & Advanced search uses 2 credits\\
Serpstack & 100 searches/month & 100 & 3.3 & Not verified & Third-party Google SERP intermediary\\
Zenserp & 50 requests/month & 50 & 1.7 & Not verified & Third-party SERP intermediary\\
\hline
\textbf{Total, Tavily basic} &  & \textbf{8,578.6} & \textbf{286.0} &  & \\
\end{tabular}
\end{center}

With Tavily advanced searches, Tavily contributes 500 searches/month rather
than 1,000:

\begin{verbatim}
8,578.6 - 500 = 8,078.6/month = 269.3/day
\end{verbatim}

A 10\% safety reserve for billing differences, failed parsing, clock boundaries,
and operational drift reduces the basic-search planning figure from 286/day to
approximately \textbf{257/day}; using a round \textbf{250/day} budget is safer.
\subsection*{Capacity by policy/risk category}
\label{sec:orga478027}

Direct/AI-search APIs only:

\begin{verbatim}
Exa 1,428.6 + Brave 1,000 + Tavily 1,000
= 3,428.6/month
= 114.3/day
\end{verbatim}

Including fixed-quota SERP intermediaries:

\begin{verbatim}
3,428.6 + Bright Data 5,000 + Serpstack 100 + Zenserp 50
= 8,578.6/month
= 286.0/day
\end{verbatim}

Bright Data supplies approximately 58\% of this broader total.  Whether the
baseline should include vendor-managed scraping of upstream search engines is
therefore a load-bearing design and policy decision.
\subsection*{Variable recurring services}
\label{sec:org96c94ca}

These services have recurring free resources, but their units cannot safely be
translated one-for-one into general-web search queries.

\begin{center}
\begin{tabular}{llll}
Service & Free resource & Best-case interpretation & Why it is not counted in the fixed total\\
\hline
ScraperAPI & 1,000 API credits/month; max 5 concurrent connections & At most 1,000 one-credit fetches/month, or 33/day & Protected domains, rendering, premium proxies, and options can consume multiple credits; it is a generic scraper rather than a structured search index\\
Apify & \char36{}5 platform credit/month & Product copy says over 1,000 Google results; roughly 100 ten-result SERP pages, approximately 3 queries/day & Actor event pricing, result count, compute, proxies, and add-ons affect cost\\
SerpAPI & Small free plan reported by secondary pages & Often reported near 250/month & Current official public page did not expose a verifiable allowance in extracted content; confirm in dashboard before counting\\
Jina Search & Keyless use and higher authenticated limits are documented & Unknown & No stable contractual free daily/monthly search allowance was found\\
\end{tabular}
\end{center}

Adding the most optimistic interpretations of ScraperAPI, Apify, and a reported
250/month SerpAPI plan to the fixed total gives roughly:

\begin{verbatim}
8,578.6 + 1,000 + 100 + 250 = 9,928.6/month
9,928.6 / 30 = 331/day
\end{verbatim}

This is an \textbf{optimistic ceiling}, not dependable capacity.
\subsection*{One-time and trial capacity}
\label{sec:orgf3de616}

\begin{center}
\begin{tabular}{llll}
Provider & One-time resource & Approximate base searches & Status/caveat\\
\hline
You.com & \char36{}100 signup credit; web search \char36{}5/1,000 & 20,000 & No card reported; does not recur\\
Exa & \char36{}20 signup credit in addition to recurring \char36{}10 & 2,857 & At \char36{}7/1,000 basic searches\\
Serper & Reported 2,500 signup credits & 2,500 & Re-verify in official dashboard before relying on it\\
SearchAPI.io & 100 get-started requests & 100 & Treat as trial unless recurrence is explicitly confirmed\\
Scrapingdog & 200 starting credits & 200 & Not shown as a recurring monthly grant\\
ScraperAPI & Up to 5,000 requests during first 7 days & Up to 5,000 & Request/credit cost can vary\\
\end{tabular}
\end{center}

\begin{verbatim}
Without the variable ScraperAPI trial: approximately 25,657 searches
Including its best-case trial: approximately 30,657 searches
At 5,000/day: approximately 5.1--6.1 days
\end{verbatim}

Signup grants must never be presented as sustainable daily capacity.
\subsection*{Existing-customer and retired products}
\label{sec:orgf688375}

\begin{itemize}
\item Google Custom Search JSON API advertises 100 free queries/day for eligible
existing customers, but it is unavailable to new customers and sunsets on
2027-01-01.  It contributes zero to a new deployment.  An already-eligible
account could temporarily add 100/day, taking the broad fixed total from
approximately 286/day to approximately 386/day.
\item Standalone Microsoft Bing Search APIs were retired on 2025-08-11 and
contribute zero.
\end{itemize}
\section*{What 5,000 calls/day requires}
\label{sec:orgb7bbb59}

With expected cache reuse:

\begin{center}
\begin{tabular}{lllll}
Cache reuse & Cache-served calls/day & Remaining live demand/day & After 286/day fixed free API capacity & Best-effort/unmetered demand\\
\hline
10\% & 500 & 4,500 & 4,500 - 286 & 4,214/day\\
40\% & 2,000 & 3,000 & 3,000 - 286 & 2,714/day\\
\end{tabular}
\end{center}

Using the safer 250/day metered budget leaves approximately 2,750--4,250 daily
queries for SearXNG, DuckDuckGo HTML, Jina, other best-effort routes, or stale
cache.

The verified fixed recurring APIs therefore cover only:

\begin{itemize}
\item Approximately 5.7\% of all 5,000 daily calls when intermediaries are included.
\item Approximately 2.3\% when only Exa, Brave, and Tavily are counted.
\item Approximately 6.4--9.5\% of the post-cache live demand with intermediaries,
depending on cache reuse.
\end{itemize}
\section*{SearXNG findings}
\label{sec:orgd6a79ee}

\subsection*{What it provides}
\label{sec:org8f09f08}

\begin{itemize}
\item A self-hosted metasearch HTTP service.
\item Query fan-out to multiple configured upstream search engines.
\item A normalized result shape over heterogeneous upstreams.
\item Engine configuration, timeout handling, and suspension/cooldown mechanisms.
\item A practical place to centralize best-effort search traffic on local hardware.
\end{itemize}
\subsection*{What it does not provide}
\label{sec:org6c8b86b}

\begin{itemize}
\item Its own general-web crawler or independent index.
\item A contractual quota from Google, Bing, DuckDuckGo, or other scraped engines.
\item Immunity from CAPTCHA challenges, rate limits, datacenter-IP blocks, or
upstream markup changes.
\item A fresh-result SLA.
\end{itemize}

SearXNG documentation explicitly discusses upstream bot classification,
CAPTCHAs, and engine suspension.  Cooldowns reduce repeated failures but do not
turn an unauthorized or blocked route into authorized capacity.

If one local SearXNG query enables four upstream engines, 5,000 Pi calls can
produce approximately 20,000 upstream requests/day.  Fan-out therefore
increases breadth and failure independence but also accelerates upstream
blocking and consumes more network work than the Pi-call count suggests.

The PDF included anecdotal DuckDuckGo blocking near 50+ queries/hour on some
shared/datacenter IPs.  This is not a guaranteed threshold; the actual limit is
unpublished, dynamic, and dependent on IP reputation and traffic shape.
\subsection*{Appropriate role}
\label{sec:org6a0c93f}

Treat SearXNG as \textbf{best-effort enrichment and overflow}, never as guaranteed free
fresh capacity.  Honor a block, CAPTCHA, or \texttt{Retry-After} by cooling down or
suspending the engine.  Do not add proxies or stealth techniques to force
continued access.
\section*{Provider and product findings}
\label{sec:org1e8901c}

\subsection*{Exa}
\label{sec:orgbf94166}

\begin{itemize}
\item Direct API and suitable result structure for Pi.
\item \char36{}20 signup credit plus \char36{}10 recurring monthly credit were documented.
\item No payment method is required for the free tier.
\item Basic search was priced at \char36{}7/1,000, yielding approximately 1,428 recurring
searches/month.
\item The terms include anti-abuse/anti-circumvention and suspension provisions.
Creating many accounts to multiply credits is therefore operationally
fragile and contrary to the design's compliance requirement.
\end{itemize}
\subsection*{Brave Search API}
\label{sec:orgd25db6e}

\begin{itemize}
\item Independent search API with a documented recurring credit.
\item \char36{}5 monthly credit at \char36{}5/1,000 searches yields 1,000/month.
\item Published capacity found during research: 50 requests/second.
\item A credit card is required even for the free plan.
\item Brave explicitly prohibits bypassing limits by creating multiple accounts.
\end{itemize}
\subsection*{Tavily}
\label{sec:orgf8240bb}

\begin{itemize}
\item 1,000 monthly credits.
\item Basic search consumes one credit; advanced consumes two.
\item Approximately 33 basic or 16 advanced searches/day on a 30-day average.
\item Development rate limit found: 100 requests/minute.
\item Existing extension support is opt-in, so it can be integrated into a
quota-aware chain without inventing a new provider implementation.
\end{itemize}
\subsection*{Bright Data SERP API}
\label{sec:org25caddb}

\begin{itemize}
\item Official docs advertise 5,000 SERP API credits every month with no credit
card.
\item Billing is per successful request; failed/errored requests are not billed in
the standard configuration.
\item Custom headers/cookies change that behavior: all attempts may then be billed.
\item Parsing, proxy management, retries, unlocking, and CAPTCHA handling are
included according to its docs.
\item It is the largest documented recurring free bucket found, but it is a
vendor-managed SERP scraper rather than first-party index access.  Its API
authorization does not automatically settle every upstream search engine's
policy.  This is a policy/legal-risk distinction, not a claim of illegality.
\end{itemize}
\subsection*{DuckDuckGo HTML}
\label{sec:orgb6ec859}

\begin{itemize}
\item Works today as the current no-key fallback.
\item It has no published programmatic quota or SLA.
\item HTML parsing can break when markup changes.
\item Rate limits and blocks are dynamic; it cannot be included in guaranteed
numeric capacity.
\end{itemize}
\subsection*{Jina Search}
\label{sec:orge0e6269}

\begin{itemize}
\item Keyless access and authenticated higher limits are documented.
\item No stable contractual free daily/monthly request allowance was found.
\item Useful only as a best-effort route until a quota is explicitly documented.
\end{itemize}
\subsection*{Serpstack, Zenserp, SerpAPI, SearchAPI, Serper, and Scrapingdog}
\label{sec:orgb30dd45}

\begin{itemize}
\item These are third-party SERP services or trials, not independent general-web
indexes.
\item Verified recurring quantities were small: Serpstack 100/month and Zenserp
50/month.
\item SerpAPI's current free quantity was not sufficiently visible in the official
extracted page to count conservatively.
\item SearchAPI, Serper, and Scrapingdog were treated as signup/trial capacity, not
recurring capacity.
\end{itemize}
\subsection*{ScraperAPI and Apify}
\label{sec:org575aa61}

\begin{itemize}
\item Both provide legitimate APIs/platforms with recurring free resources.
\item They are generic scraping infrastructure or actors, so one credit, one
result, one SERP page, and one user query are different units.
\item They can be optional overflow experiments, but the router must learn actual
observed credit cost before assigning a hard daily budget.
\end{itemize}
\section*{Consumer subscriptions and hosted agents}
\label{sec:org22e9aaa}

The following were considered and rejected as raw capacity sources:

\begin{itemize}
\item ChatGPT subscriptions.
\item Gemini consumer subscriptions.
\item Perplexity subscriptions.
\item Claude subscriptions.
\item OpenCode Go.
\item Kagi and Perplexity paid APIs under a hard-\char36{}0 budget.
\end{itemize}

Reasons:

\begin{itemize}
\item Consumer plans provide hosted interactive research, not 5,000 raw
programmatic search calls/day for Pi.
\item Automating a consumer UI would be fragile and commonly disallowed.
\item A hosted research agent would control query planning and result handling,
violating the requirement that Pi orchestrate every query and receive every
result.
\item OpenCode Go covers model inference rather than increased web-search capacity.
\item OpenCode's unauthenticated Exa MCP route has no appropriate published quota
or SLA and is not a basis for production planning.
\end{itemize}
\section*{Compliance and routing rules}
\label{sec:orgb9e99f9}

\subsection*{Allowed}
\label{sec:org86a98b1}

\begin{itemize}
\item One legitimate account per independent provider.
\item Weighted or quota-aware routing across providers.
\item Pacing requests below each published QPS/RPM limit.
\item Tracking monthly credits locally and stopping before exhaustion.
\item Honoring \texttt{429}, \texttt{Retry-After}, CAPTCHA, and explicit blocks.
\item Caching, in-flight deduplication, stale fallback, and partial responses.
\item Using a provider's documented API according to its own terms.
\end{itemize}
\subsection*{Excluded}
\label{sec:org8b7300d}

\begin{itemize}
\item Waiting for one account to be blocked and continuing through another account
at the same provider.
\item Multiple accounts, identities, or organizations created to multiply a free
tier.
\item IP/proxy rotation intended to defeat an upstream limit.
\item Browser fingerprint spoofing or stealth plugins for this feature.
\item CAPTCHA solving or bypass.
\item Describing an intermediary API as first-party authorization from the upstream
engine.
\end{itemize}

Provider rotation increases resilience only when each route is independently
authorized.  It does not convert scraping or quota bypass into authorized use.
\section*{Approaches considered}
\label{sec:orga131658}

\subsection*{Approach 1: API-only free-tier router}
\label{sec:org3f68652}

Use Exa, Brave, Tavily, and small documented free APIs with persistent quota
tracking.

Advantages:

\begin{itemize}
\item Clear result contracts.
\item Easier accounting and circuit breaking.
\item Better policy position than direct scraping.
\end{itemize}

Disadvantages:

\begin{itemize}
\item Only approximately 114/day from direct/AI-search APIs.
\item Only approximately 286/day even after adding fixed SERP intermediaries.
\item Cannot approach 5,000 fresh searches/day.
\end{itemize}

Conclusion: useful as the high-quality tier, insufficient by itself.
\subsection*{Approach 2: resilient local search gateway (recommended baseline)}
\label{sec:org661f093}

Combine persistent cache, quota-aware authorized APIs, self-hosted SearXNG,
best-effort direct routes, stale fallback, and explicit metadata.

Advantages:

\begin{itemize}
\item Pi retains full orchestration.
\item Can target high non-empty-response availability.
\item Uses scarce authorized API credits where they add the most value.
\item Persistent state survives Pi processes and reports.
\item Failures are explicit and measurable.
\end{itemize}

Disadvantages:

\begin{itemize}
\item Cannot guarantee freshness.
\item SearXNG/upstream parsing remains operationally fragile.
\item 99\% availability depends materially on cached/stale/partial acceptance.
\item More components and policy distinctions must be maintained.
\end{itemize}

Conclusion: best fit under the clarified requirements and \char36{}0 constraint.
\subsection*{Approach 3: local crawler/index}
\label{sec:orga792812}

Build or import a local searchable web corpus.

Advantages:

\begin{itemize}
\item Predictable local availability after indexing.
\item No per-query provider quota for already indexed documents.
\end{itemize}

Disadvantages:

\begin{itemize}
\item A broad fresh web index has major storage, bandwidth, crawling, ranking, and
compliance costs.
\item Poor fit for a hard-\char36{}0 general-web freshness target.
\item Estimated query reuse of 10--40\% may not justify the initial complexity.
\end{itemize}

Conclusion: defer.  Reconsider only if observed query/domain reuse supports a
narrow, valuable local corpus.
\section*{Recommended design baseline}
\label{sec:org3791195}

This remains a proposal pending explicit approval.
\subsection*{Components}
\label{sec:orgdc96810}

\begin{enumerate}
\item \textbf{Gateway inside or adjacent to the web-search extension}
\begin{itemize}
\item Receives normalized \texttt{web\_lookup} requests.
\item Owns routing, deadlines, merge behavior, and response metadata.
\end{itemize}
\item \textbf{Persistent SQLite cache}
\begin{itemize}
\item Shared across Pi sessions and processes.
\item Stores normalized query/parameters, results, timestamps, provider,
freshness class, and quality/failure metadata.
\end{itemize}
\item \textbf{In-flight request coalescing}
\begin{itemize}
\item Concurrent equivalent queries share one upstream request.
\end{itemize}
\item \textbf{Provider adapters}
\begin{itemize}
\item Existing Exa, Tavily, and DuckDuckGo adapters.
\item Optional Brave, Bright Data, SearXNG, and other approved adapters.
\end{itemize}
\item \textbf{Quota ledger and rate limiter}
\begin{itemize}
\item Persistent monthly counters plus in-memory token buckets for burst limits.
\end{itemize}
\item \textbf{Engine health manager}
\begin{itemize}
\item Consecutive-failure counts, cooldowns, circuit breakers, and recovery
probes.
\end{itemize}
\item \textbf{Merger/deduplicator}
\begin{itemize}
\item Canonicalizes URLs, removes duplicates, preserves provenance, and merges
non-identical useful snippets.
\end{itemize}
\item \textbf{Metrics and operator visibility}
\begin{itemize}
\item Measures status, freshness, provider use, quota, latency, and failures.
\end{itemize}
\end{enumerate}
\subsection*{Routing sequence}
\label{sec:org35a29f4}

\begin{verbatim}
normalize query and parameters
        |
        v
fresh persistent cache hit? -- yes --> return status=fresh-cache
        |
        no
        v
coalesce with equivalent in-flight query if present
        |
        v
choose healthy provider with quota and rate tokens
        |
        +--> authorized/direct APIs
        +--> approved SERP APIs
        +--> self-hosted SearXNG / best-effort routes
        |
        v
merge + URL-deduplicate + persist
        |
        +--> non-empty live result: return status=fresh or partial
        |
        +--> live routes unavailable/empty
                    |
                    v
             stale cache hit? -- yes --> return status=stale
                    |
                    no
                    v
             explicit partial/unavailable response
\end{verbatim}

The exact ordering within the live tier should be configurable and may depend
on query type, remaining monthly budget, historical quality, and latency.
\subsection*{Query normalization and cache keys}
\label{sec:org043906f}

At minimum, keys need to include:

\begin{itemize}
\item Whitespace and Unicode-normalized query text.
\item Case normalization where semantics permit it.
\item Requested engine or automatic routing mode.
\item Result limit.
\item Locale, country, language, safe-search, and time filters.
\item Any search mode that materially changes results.
\item A cache schema/version field so parser or ranking changes can invalidate old
entries safely.
\end{itemize}

Do not normalize away quoted phrases, operators, punctuation, or filters that
change search meaning.
\subsection*{Freshness model}
\label{sec:orgfc99731}

Response states should distinguish at least:

\begin{itemize}
\item \texttt{fresh}: produced by a live upstream call within the request.
\item \texttt{fresh-cache} or \texttt{cached}: cache entry within configured fresh TTL.
\item \texttt{stale}: cache entry beyond fresh TTL but within allowed stale retention.
\item \texttt{partial}: non-empty result set with one or more providers failed, timed out,
or were skipped.
\item \texttt{unavailable}: no non-empty live or cache result; should not be counted as an
availability success.
\end{itemize}

Exact fresh TTLs and maximum stale ages were not approved during early
research.  They should be configurable by query/result class and selected from
observed reuse and freshness needs rather than invented globally.

Recommended metadata fields:

\begin{verbatim}
{
  "status": "fresh|cached|stale|partial|unavailable",
  "query": "normalized query",
  "results": [],
  "providersAttempted": [],
  "providersSucceeded": [],
  "partialFailures": [],
  "fetchedAt": "timestamp",
  "cacheAgeSeconds": 0,
  "freshnessTtlSeconds": 0,
  "servedFromCache": false,
  "quotaSnapshot": {},
  "requestId": "..."
}
\end{verbatim}

The public shape may need to preserve backward compatibility with the current
\texttt{web\_lookup} tool while exposing additional metadata in a concise text or
structured envelope.
\subsection*{Quota-aware scheduling}
\label{sec:orgde30be6}

Do not implement "try providers until one emits 429" as normal scheduling.
Maintain local counters and reserve capacity across the remaining billing
period.

For a monthly allowance:

\begin{verbatim}
safe_remaining = documented_remaining - reserve
pace_per_day = safe_remaining / days_remaining_in_provider_cycle
\end{verbatim}

A weighted distribution over the 286/day broad total is approximately:

\begin{center}
\begin{tabular}{ll}
Provider & Approximate share\\
\hline
Bright Data & 58.3\%\\
Exa & 16.7\%\\
Brave & 11.7\%\\
Tavily basic & 11.7\%\\
Serpstack & 1.2\%\\
Zenserp & 0.6\%\\
\end{tabular}
\end{center}

Use those figures only after each provider is approved and configured.  A safer
operational allocation is approximately:

\begin{itemize}
\item Bright Data: 150/day.
\item Exa: 42/day.
\item Brave: 30/day.
\item Tavily basic: 30/day.
\item Serpstack: 3/day.
\item Zenserp: 1/day.
\item Total: approximately 256/day, with monthly reserve.
\end{itemize}

The scheduler should also support:

\begin{itemize}
\item Per-second/per-minute token buckets.
\item Provider billing-cycle boundaries rather than assuming every month is exactly
30 days.
\item Atomic quota updates across concurrent Pi processes.
\item Manual quota correction after dashboard reconciliation.
\item Hard zero-cost behavior: never fall through to paid overage.
\item Provider disable switches.
\item Priority preservation for high-value or freshness-sensitive requests if the
user later defines such classes.
\end{itemize}
\subsection*{Cooldowns and circuit breakers}
\label{sec:orgc663866}

Suggested state transitions:

\begin{itemize}
\item Timeout/transient 5xx: short exponential backoff with jitter.
\item \texttt{429} with \texttt{Retry-After}: disable until the indicated time.
\item Repeated failures or parser-empty anomalies: open circuit for a longer
cooldown.
\item CAPTCHA/block page: suspend route; do not retry aggressively.
\item Quota exhausted: disable until the next verified billing reset.
\item Half-open recovery: one bounded probe, then close or reopen.
\end{itemize}

Track failures separately by engine and failure class.  An empty result can be
legitimate for a narrow query, so parser/block detection must not equate every
empty result with upstream failure.
\subsection*{Result merging}
\label{sec:org60073fe}

\begin{itemize}
\item Normalize URLs by removing known tracking parameters and fragments where
safe.
\item Preserve canonical URL, original URL, provider, title, snippet, rank, and
retrieval time.
\item Deduplicate exact/canonical URLs across providers.
\item Prefer more informative titles/snippets without discarding provenance.
\item Do not pretend independently returned duplicates are independent evidence.
\item Record which providers contributed each merged result.
\item Keep provider-native ordering available for debugging.
\end{itemize}
\subsection*{Observability}
\label{sec:org89749b5}

At minimum collect:

\begin{itemize}
\item Total calls and non-empty response rate.
\item Counts by \texttt{fresh}, \texttt{cached}, \texttt{stale}, \texttt{partial}, and \texttt{unavailable}.
\item Cache hit rate and cache-age distributions.
\item Per-provider attempts, successes, empty responses, timeouts, blocks, 429s,
latency, and result count.
\item Quota remaining and projected exhaustion date.
\item In-flight deduplication savings.
\item URL overlap among providers.
\item Percentage of requests requiring stale fallback.
\item Percentage of partial responses and their failed providers.
\end{itemize}

These metrics are necessary to distinguish "the tool answered" from "the web
was searched freshly and comprehensively."
\section*{Service objectives and honest claims}
\label{sec:org2edb433}

\subsection*{Feasible target}
\label{sec:orgd617df4}

A target such as the following is plausible only after measurement:

\begin{itemize}
\item 99\% of calls return a non-empty response from live results, fresh cache, stale
cache, or an explicitly partial non-empty result.
\item Every response includes freshness and degradation metadata.
\item Zero paid overage.
\item No retry behavior that circumvents provider controls.
\end{itemize}
\subsection*{Not supportable under these constraints}
\label{sec:org05eb1b4}

\begin{itemize}
\item 99\% fresh general-web result availability at 5,000/day.
\item 99\% comprehensive coverage of the public web.
\item A guarantee that SearXNG's famous-engine backends remain usable.
\item A claim that 5,000 successful local SearXNG HTTP responses equal 5,000 fresh,
independent search-index queries.
\item A fixed number for Jina, DuckDuckGo HTML, or public SearXNG capacity without a
published quota/SLA.
\end{itemize}
\subsection*{Availability denominator}
\label{sec:orgee8b813}

Before implementation, define precisely:

\begin{itemize}
\item Whether a partial response with one result counts.
\item Whether an informative \texttt{unavailable} response counts only as tool uptime or
as search success (recommended: tool uptime yes, search success no).
\item Maximum acceptable stale age by use case.
\item Whether duplicate-only results count.
\item Whether a cache hit for exactly normalized parameters is required.
\end{itemize}

Report at least two objectives separately:

\begin{enumerate}
\item \textbf{Tool availability}: the function returns a valid response envelope.
\item \textbf{Search-result availability}: the response contains one or more results.
\end{enumerate}

Also report live/fresh coverage separately so stale fallback cannot hide an
upstream outage.
\section*{Rejected strategies and rationale}
\label{sec:org926003a}

\subsection*{Many Exa or Brave accounts}
\label{sec:org0feeb72}

Rejected because it is quota multiplication rather than independent-provider
resilience.  Brave explicitly prohibits multiple-account limit bypass; Exa has
anti-circumvention and suspension language.  It also creates key-management,
identity, reset, and sudden-ban failure modes.
\subsection*{Playwright/stealth automation of search engines}
\label{sec:org9d39a15}

Rejected because browser mimicry does not create stable authorization or quota.
CAPTCHA and block thresholds are dynamic, browser maintenance is costly, and
forcing continued access would violate the no-circumvention requirement.
\subsection*{Proxy rotation}
\label{sec:org1df2b14}

Rejected under the \char36{}0 and compliance constraints.  Free proxies are unreliable
and unsafe; paid proxies violate the budget; rotation intended to evade blocks
is out of scope regardless of cost.
\subsection*{Consumer research subscriptions as backends}
\label{sec:orgf5a8af9}

Rejected because they do not grant Pi a raw high-volume search API.  They also
move query planning and evidence handling into another hosted agent.
\subsection*{Claiming burst aggregation solves monthly capacity}
\label{sec:org3238d64}

Rejected mathematically.  Fifty requests/second is irrelevant once a provider's
1,000 monthly requests are consumed.  Burst limits control when requests may be
sent; credits control how many can be sent sustainably.
\section*{Open decisions before implementation}
\label{sec:org1d29c4e}

\begin{itemize}
\item Approve or reject Approach 2 as the design baseline.
\item Decide whether Bright Data and other third-party SERP scrapers are acceptable
policy dependencies.  Without them, recurring numeric capacity falls from
approximately 286/day to approximately 114/day.
\item Decide which providers may require a payment card even though configured
spend remains zero (notably Brave).
\item Define fresh TTLs and maximum stale retention by query class.
\item Define the exact success denominator for the 99\% objective.
\item Decide whether live providers race in parallel, run sequentially, or use a
bounded hybrid.  Parallel racing improves latency but consumes multiple
provider credits for one user query.
\item Choose whether provider results are first-success or merged.  Merging improves
coverage but multiplies quota use.
\item Choose SQLite location, locking strategy, retention size, and privacy policy.
\item Define cache handling for potentially sensitive queries.
\item Decide whether explicit \texttt{engine} requests may fall back to other engines or
must fail closed.
\item Reconfirm every provider's quota, reset behavior, overage behavior, and terms
immediately before implementation because free plans change frequently.
\item Determine whether the extension should run SearXNG itself, call an externally
managed local instance, or merely support an optional URL.
\item Establish a test method using mocked provider responses; live tests must not
consume uncontrolled quota.
\end{itemize}
\section*{Likely implementation impact}
\label{sec:org6ea1e47}

If the design is approved, likely touched or new areas include:

\begin{itemize}
\item \texttt{extensions/web-search/search.ts}: replace first-success-only chain with the
approved cache/router behavior.
\item \texttt{extensions/web-search/types.ts}: add provider health, provenance, cache, and
response-status types.
\item \texttt{extensions/web-search/index.ts}: configuration/CLI options and result
rendering while preserving tool compatibility.
\item \texttt{extensions/web-search/engines/}: add approved Brave, Bright Data, and SearXNG
adapters.
\item New cache/quota/health modules under \texttt{extensions/web-search/}.
\item \texttt{tests/web-search.test.ts}: deterministic cache, quota, cooldown, fallback,
merge, and metadata tests.
\item \texttt{skills/web-search/SKILL.md}: teach the model how to interpret fresh, cached,
stale, partial, and unavailable responses.
\end{itemize}

This is only a scope sketch.  A written implementation plan and test-first
workflow should follow design approval.
\section*{Research quality, contradictions, and gaps}
\label{sec:orgaf8e80d}

\begin{itemize}
\item Provider pricing and free plans are volatile.  The values in this document
are snapshots from official documentation reviewed during this research.
\item "All available providers" cannot literally be exhaustive: scraper
marketplaces, promotional credits, and small SERP vendors appear and vanish
frequently.
\item Search-result snippets sometimes disagreed with current official pages.  The
conservative calculation counts only quantities that could be tied to
official provider material, except where explicitly labeled reported or
provisional.
\item Bright Data's documentation maps its 5,000 monthly credits to successful SERP
requests and says unlocking/CAPTCHA handling is included.  This supports the
capacity calculation but also reinforces that it is a managed scraper, not
first-party search-index authorization.
\item ScraperAPI and Apify have real recurring free resources, but a stable
query-equivalent requires actual configuration and billing telemetry.
\item SerpAPI's recurring free quantity remained insufficiently verified from the
official extracted page and is excluded from the conservative total.
\item Jina's keyless access lacks a stable published hard allowance and is excluded
from numeric capacity.
\item SearXNG and DuckDuckGo can add substantial best-effort volume, but assigning a
hard daily number would be false precision.
\item A \char36{}0 system cannot make a contractual availability guarantee when its bulk
live capacity comes from services with no SLA.
\end{itemize}
\section*{Primary references}
\label{sec:org8df449a}

Official and primary documentation used or identified during the research:

\begin{itemize}
\item Bright Data SERP quickstart and free tier:
\url{https://docs.brightdata.com/scraping-automation/serp-api/quickstart}
\item Bright Data SERP pricing and successful-request billing:
\url{https://docs.brightdata.com/scraping-automation/serp-api/pricing-and-billing}
\item Exa billing/pricing documentation:
\url{https://exa.ai/docs/reference/billing}
\item Exa terms:
\url{https://exa.ai/terms}
\item Brave Search API pricing:
\url{https://api-dashboard.search.brave.com/documentation/pricing}
\item Brave Search API terms:
\url{https://brave.com/search/api/terms-of-service/}
\item Tavily API credits:
\url{https://docs.tavily.com/documentation/api-credits}
\item Serpstack pricing:
\url{https://serpstack.com/pricing}
\item Zenserp pricing:
\url{https://zenserp.com/pricing-plans/}
\item ScraperAPI plans and billing:
\url{https://docs.scraperapi.com/resources/faq/plans-and-billing}
\item ScraperAPI credit/request costs:
\url{https://docs.scraperapi.com/getting-started/quick-start/credits-and-requests-costs}
\item Apify pricing:
\url{https://apify.com/pricing}
\item Apify Google Search Scraper:
\url{https://apify.com/apify/google-search-scraper}
\item You.com billing:
\url{https://you.com/docs/administration/billing}
\item Google Custom Search JSON API overview:
\url{https://developers.google.com/custom-search/v1/overview}
\item Microsoft Bing Search API retirement:
\url{https://learn.microsoft.com/en-us/lifecycle/announcements/bing-search-api-retirement}
\item SearXNG administrator documentation:
\url{https://docs.searxng.org/admin/index.html}
\item SearXNG bot/CAPTCHA documentation:
\url{https://docs.searxng.org/admin/answer-captcha.html}
\item Jina Reader/Search documentation:
\url{https://jina.ai/reader/}
\item Scrapingdog Google SERP API:
\url{https://www.scrapingdog.com/google-serp-api/}
\item Scrapingdog pricing:
\url{https://www.scrapingdog.com/pricing/}
\end{itemize}
\section*{Current recommendation}
\label{sec:org8706914}

Adopt Approach 2 as the design baseline only if the availability goal is stated
honestly as high \textbf{non-empty response availability}, with separate reporting for
fresh/live coverage.  Use a persistent cache and quota-aware authorized APIs as
the reliable core, SearXNG and other no-SLA routes as best-effort enrichment,
and stale/partial responses as explicit degradation modes.

Do not claim that this architecture produces 5,000 fresh searches/day at \char36{}0.
It produces up to 5,000 Pi-controlled lookup attempts/day and seeks to answer
most of them through a measured combination of live search, cache, stale data,
and partial results.
\end{document}
